\section{Results at a glance}
\label{sec:results-map}

Under its hypotheses, Theorem~\ref{thm:main} classifies every coset of
weight at least \(r-1\).
Its conclusion is independent of the placement of the deleted set \(A\) and
of all nonzero GRS column multipliers.

\begin{table}[H]
\centering
\small
\caption{The high-weight cosets in the classification range of
Theorem~\ref{thm:main}: \(q\geq Q^*_{r,s}\), and \(p\) odd or \(p=2\)
with \(r\geq8\).  Counts are projective syndrome directions;
literal nonzero coset counts are larger by a factor of \(q-1\).}
\label{tab:headline-shells}
\setlength{\tabcolsep}{5pt}
\begin{tabular}{p{0.12\textwidth}p{0.12\textwidth}p{0.25\textwidth}
                p{0.38\textwidth}}
\toprule
Support & Radius & Weight \(r\) & Weight \(r-1\)\\
\midrule
\(s=0\) & \(r-1\) & none & tangents and conjugate secants;
count \(q(q+1)^2/2\)\\
\(s>0\) & \(r\) & omitted curve points; count \(s\) & tangents,
conjugate secants, and interiors of split secants incident with \(A\);
count \(q(q+1)^2/2+(q-1)s(2q+1-s)/2\)\\
\bottomrule
\end{tabular}
\end{table}

Appending a syndrome direction as a parity-check column raises its coset
weight by one.  The preceding table therefore gives all one-column
extensions with Singleton defect at most one:
\[
\begin{array}{c|c|c}
d_S(f)&d(\ker[H_S\mid f])&\text{extension class}\\ \hline
r&r+1&\text{MDS}\\
r-1&r&\text{NMDS}\\
\leq r-2&\leq r-1&\text{defect at least two}.
\end{array}
\]
Theorem~\ref{thm:family-aggregate-nmds} then counts the minimum supports in
each NMDS family and determines its complete aggregate weight enumerator.

In the two remaining binary cases the general conclusion is
\[
 d_S(f)\geq r-1
 \quad\Longrightarrow\quad
 f\in\cP_r\cup\cM^{\max}_{r,2},\qquad r\in\{6,7\},
\]
and, for full support, Theorems~\ref{thm:r6} and~\ref{thm:r7} decide the
arithmetic of that carrier over every prime power \(q\geq7\).  The fixed R5--R7 classifications
sharpen the uniform theorem in small fields but are not inputs to it.

\subsection{Why the classification closes}

The proof has three steps.  The Hankel kernel \(W_f\) records the
degree-\((r-2)\) binary forms whose roots can represent \(f\): a squarefree
product of rational linear factors supported on \(S\) is exactly an
\((r-2)\)-column representation.  Choosing one factor contracts both the
syndrome and the witness problem by one degree; choosing \(r-5\) factors
leaves a redundancy-five pencil of binary cubics.  The recursive carrier
theorem identifies when every such contraction is trapped.  Everywhere else,
one degree-six selector chooses the factors simultaneously, and the exact R5
pencil count supplies a split squarefree cubic avoiding those factors and \(A\).

\begin{figure}[H]
\centering
\[
\begin{array}{c}
f\notin
\underbrace{\cP_r}_{\text{persistent rank-two locus}}
\ \bigl(\cup\ \cM^{\max}_{r,2}\text{ at }(r,p)\in\{(6,2),(7,2)\}\bigr)\\[1mm]
\Downarrow\quad\text{simultaneous selection and contraction}\\[1mm]
\substack{\text{one split squarefree cubic in the terminal Hankel pencil,}\\
          \text{avoiding the markers and }A}\\[1mm]
\Downarrow\quad\text{lift through }r-5\text{ distinct markers in }S\\[1mm]
W_f\text{ contains a split squarefree form supported on }S\\[1mm]
\Updownarrow\quad\text{Hankel criterion}\\[1mm]
d_S(f)\le r-2.
\end{array}
\]
\caption{The proof route.  The carrier theorem identifies the only possible
obstruction; simultaneous contraction constructs a locator everywhere else.}
\label{fig:headline-proof-route}
\end{figure}

For the fixed-level bounded fields, exact certificates enumerate projective
syndromes modulo the stated group and record
canonical representatives, stabilizers, orbit sizes, and split-member data
relative to the stated evaluation set,
and check that the orbit sizes exhaust the ambient projective space.  A
separate replay recomputes those invariants from the records.  These finite
steps complete the split-free geometric classification only in the named fields;
covering-radius promotion to deep holes remains a separate cited input.

A first-pass reader can now continue with the coding dictionary in
Section~\ref{sec:dictionary}, the carrier and escape theorems in
Sections~\ref{sec:recursive-carriers}--\ref{sec:simultaneous-contraction},
and their coding consequences in Section~\ref{sec:high-weight-cosets}.
Section~\ref{sec:r5} then proves the terminal pencil input used by that
spine.  The detailed fixed-level polar packages, exceptional orbit tables,
modular calculations, and verification records are collected in the
appendices.
